% =====================================================================
% A Newton-polygon formal-level baseline for the Birkhoff-Trjitzinsky
% exponent on the SIARC polynomial continued-fraction Wallis stratum
%
% Standalone SIARC short note (Zenodo cite-target; new concept DOI).
%
% Provenance: this note consolidates the result of T1-BIRKHOFF-PHASE2-
% LIFT-LOWER (verdict UPGRADE_PARTIAL_FORMAL_LEVEL; bridge commit
% 37c939f, 2026-05-04). It introduces no new numerical pipeline:
% every numerical fact is carried over from the Phase 2 bridge
% session and re-anchored in this note's own claims.jsonl.
%
% Date: 2026-05-06
% Filename convention: bt_baseline_note (presentation-grade)
% =====================================================================

\documentclass[11pt,reqno]{amsart}

\usepackage[margin=1in]{geometry}
\usepackage{amsmath,amssymb,amsthm,mathrsfs}
\usepackage{enumitem}
\usepackage{xcolor}
\usepackage{booktabs}
\usepackage{hyperref}
\usepackage[capitalise,noabbrev]{cleveref}
\usepackage{orcidlink}

\definecolor{darklink}{rgb}{0.0,0.2,0.5}
\hypersetup{
  colorlinks=true,
  linkcolor=darklink,
  citecolor=darklink,
  urlcolor=darklink,
  pdftitle={A Newton-polygon formal-level baseline for the
            Birkhoff-Trjitzinsky exponent on the SIARC polynomial
            continued-fraction Wallis stratum (v1.0)},
  pdfauthor={Papanokechi},
  pdfkeywords={Newton polygon, Birkhoff-Trjitzinsky, polynomial
               continued fractions, SIARC PCF stratum, asymptotic
               exponent, Wimp-Zeilberger ansatz}
}

\theoremstyle{plain}
\newtheorem{theorem}{Theorem}[section]
\newtheorem{lemma}[theorem]{Lemma}
\newtheorem{proposition}[theorem]{Proposition}
\newtheorem{corollary}[theorem]{Corollary}
\newtheorem{conjecture}[theorem]{Conjecture}

\theoremstyle{definition}
\newtheorem{definition}[theorem]{Definition}
\newtheorem{example}[theorem]{Example}
\newtheorem{problem}[theorem]{Open problem}

\theoremstyle{remark}
\newtheorem{remark}[theorem]{Remark}

\newcommand{\PCF}{\mathrm{PCF}}
\newcommand{\dps}{\mathrm{dps}}
\newcommand{\Anaive}{A_{\mathrm{naive}}}
\newcommand{\Afit}{A_{\mathrm{fit}}}
\newcommand{\dega}{\deg a}
\newcommand{\degb}{\deg b}

\def\version{1.0}

\title[A Newton-polygon formal-level baseline]
      {A Newton-polygon formal-level baseline for the
       Birkhoff--Trjitzinsky exponent on the SIARC polynomial
       continued-fraction Wallis stratum}

\author{Papanokechi}
\address{Independent researcher, Yokohama, Japan. \\
ORCID: \href{https://orcid.org/0009-0000-6192-8273}{0009-0000-6192-8273}.}
\thanks{ORCID:
\href{https://orcid.org/0009-0000-6192-8273}{0009-0000-6192-8273}.
The author received no external funding for this work. The symbolic
derivations carried by this note were performed under the
AEAL/SIARC multi-agent computational methodology; AI assistants
(GitHub~Copilot powered by Anthropic Claude Opus~4.7, and
Anthropic Claude directly) were used for code generation, numerical
computations, and manuscript drafting. After using these tools, the
author reviewed and edited the content as needed and takes full
responsibility for the content of this article. The AI assistants
are not authors. Session-level computation logs and prompts are
archived on the SIARC bridge under
\texttt{sessions/2026-05-03/T1-BIRKHOFF-PHASE2-LIFT-LOWER/}
(commit \texttt{37c939f}) and
\texttt{sessions/2026-05-06/T1-PHASE2-BASELINE-NOTE/} (this
version), both available in line with the AEAL reproducibility
requirement.}

\date{2026-05-06 (v\version)}

\begin{document}

\begin{abstract}
For the SIARC polynomial continued-fraction (PCF) Wallis recurrence
$p_n = b_n\,p_{n-1} + a_n\,p_{n-2}$ with polynomial coefficients
$b_n$ of degree $d \ge 2$ and $a_n$ of convention-dependent degree,
we record the formal-level baseline obtained from the standard
Wimp--Zeilberger normal-case ansatz
$y(n) = \Gamma(n)^{\mu}\,\gamma^{n}\,n^{\rho}\,S(1/n)$:
the formal asymptotic exponent
$\Anaive = \mu_{\mathrm{dom}} - \mu_{\mathrm{sub}}$ takes a value
in $\{d-1,\,d,\,d+1\}$ across the three coefficient conventions
(\emph{$\alpha$}-direction, symmetric, \emph{$\delta$}-direction)
uniformly for $d \in [2,8]$ (\Cref{thm:baseline}). The
$d=2$, $\alpha$-direction prediction $\Anaive = 3$ is verified
at $80$ algebraic digits against PCF-1 v1.3 \S6 Theorem 5 lower
branch (QL01--QL26)
\cite{papanokechi_pcf1_v13}\footnote{SIARC Zenodo deposit;
internal review only, not peer-reviewed.}. The empirical record of
PCF-2 v1.3 R1.1 + R1.3 + Q1
(\cite{papanokechi_pcf2_v13}\footnote{SIARC Zenodo deposit; internal
review only, not peer-reviewed.}: $50/50$ cubic representatives at
$d=3$ and $60/60$ quartic representatives at $d=4$ exhibiting
$\Afit \approx 2d$) lies strictly outside the formal baseline
at $d \ge 3$; \Cref{prop:gap-framing} frames this gap by
identifying two non-mutually-exclusive structural mechanisms
consistent with it. The note remains purely consolidatory
(no new numerical pipeline is run) and is the SIARC counterpart
of PCF-1 v1.3 \S6 Theorem 5 \emph{lower branch}: a positive
formal-level result in its own right, citable independently of
any subsequent lift to $A = 2d$
(Conjecture B4 of \cite{papanokechi_pcf2_v13}).
\end{abstract}

\maketitle

\section{Setup and main results}
\label{sec:setup}

\subsection{The SIARC PCF Wallis recurrence}
Throughout this note we work with a polynomial continued-fraction
(PCF) Wallis recurrence in the SIARC convention
\begin{equation}
\label{eq:wallis}
  p_n \;=\; b_n\,p_{n-1} \,+\, a_n\,p_{n-2},
  \qquad n \ge 2,
\end{equation}
with the partial denominator polynomial
$b_n = b(n) = c_b\,n^{d} + (\text{lower order in }n)$
of degree $d \ge 2$ and leading coefficient $c_b \ne 0$, and the
partial numerator polynomial $a_n = a(n)$ a polynomial of degree
$\dega$ that depends on the SIARC \emph{coefficient convention}
in use. We adopt the three conventions standard in the
SIARC stack:

\begin{itemize}[leftmargin=2.2em]
\item \textbf{$\alpha$-direction.}
  $\dega = d - 1$.
  This is the convention of PCF-1 v1.3 \S6 (the Wallis-product
  family with one fewer numerator degree than denominator)
  and of the QL01--QL26 family
  \cite{papanokechi_pcf1_v13}.
\item \textbf{Symmetric.}
  $\dega = d$.
  This is the convention used for several intermediate
  asymptotic-classification arguments in the Channel Theory
  outline \cite{siarc_channel_theory_v13}.
\item \textbf{$\delta$-direction.}
  $\dega = d + 1$.
  This is the convention of the V$_{\mathrm{quad}}$ class label
  and is invoked in the gap-proposition framing of the
  Birkhoff--Trjitzinsky lift \cite{siarc_t1_phase2}.
\end{itemize}

The leading coefficient of $a_n$ is denoted $c_a \ne 0$ in all
three conventions.

\subsection{The Wimp--Zeilberger normal-case ansatz}
\label{sec:wz-ansatz-intro}
The standard formal-asymptotic ansatz for the second-order linear
difference recurrence \eqref{eq:wallis} at the irregular singular
point $n = \infty$ is the Wimp--Zeilberger normal-case ansatz
\cite[Ch.~6]{wimp1984computation}, in the form
\begin{equation}
\label{eq:wz-ansatz}
  y(n) \;=\; \Gamma(n)^{\mu}\;\gamma^{n}\;n^{\rho}\;S(1/n),
  \qquad
  S(t) \;=\; 1 \,+\, s_1\,t \,+\, s_2\,t^{2} \,+\, \dots,
\end{equation}
a formal series with $\mu, \gamma, \rho \in \mathbb{C}$ and
$S(t) \in \mathbb{C}[[t]]$. The objects $\mu$ and $\gamma$ are
determined by the Newton-polygon balance of the leading
coefficients of the recurrence; $\rho$ is the secondary
indicial exponent; and $S(1/n)$ encodes the Birkhoff
recursion for the lower-order corrections.

In this note, the formal asymptotic exponent of the recurrence
\eqref{eq:wallis} at the irregular singular point $n = \infty$
is the Birkhoff--Trjitzinsky difference
\begin{equation}
\label{eq:Anaive-def}
  \Anaive \;:=\; \mu_{\mathrm{dom}} \,-\, \mu_{\mathrm{sub}},
\end{equation}
where $\mu_{\mathrm{dom}}$ and $\mu_{\mathrm{sub}}$ are the
$\mu$-values of, respectively, the dominant and subdominant
formal solutions of \eqref{eq:wallis} produced by the
ansatz \eqref{eq:wz-ansatz}. The subscript \emph{naive}
records that \eqref{eq:Anaive-def} is the formal-level
exponent obtained from the normal-case ansatz; the relation
of $\Anaive$ to the empirical exponent $\Afit$ measured in
high-precision tail-window fits of $\log|p_n / p_{n-1}|$ on
the SIARC stratum at $d \ge 3$ is the substance of
\Cref{sec:gap}.

\subsection{Main results}
\label{sec:main-results}

The two main results of this note are the following.

\begin{theorem}[Formal-level baseline]
\label{thm:baseline}
For the SIARC PCF Wallis recurrence \eqref{eq:wallis} with
polynomial coefficient $b_n$ of degree $d \ge 2$ (leading
coefficient $c_b \ne 0$) and polynomial coefficient $a_n$ of
convention-dependent degree
$\dega \in \{d-1,\,d,\,d+1\}$
(leading coefficient $c_a \ne 0$), the Wimp--Zeilberger
normal-case ansatz \eqref{eq:wz-ansatz} admits exactly two
formal-balance solutions $y_{\mathrm{dom}}$ and
$y_{\mathrm{sub}}$, with
\begin{equation}
\label{eq:mu-pair}
  \mu_{\mathrm{dom}} \;=\; \degb \;=\; d,
  \qquad
  \mu_{\mathrm{sub}} \;=\; \dega \,-\, \degb \;=\; \dega - d,
\end{equation}
so that the formal asymptotic exponent
\eqref{eq:Anaive-def} satisfies
\begin{equation}
\label{eq:Anaive-band}
  \Anaive \;=\; 2\degb - \dega \;=\; 2d - \dega
  \;\in\; \{\,d-1,\,d,\,d+1\,\}
\end{equation}
across the three SIARC conventions, uniformly for
$d \in [2, 8]$.
\end{theorem}

\begin{proposition}[Gap-framing]
\label{prop:gap-framing}
The empirical record of PCF-2 v1.3 R1.1 + R1.3 + Q1
\cite{papanokechi_pcf2_v13} -- $50/50$ cubic representatives
at $d=3$ with $\Afit \approx 6 = 2d$, and $60/60$ quartic
representatives at $d=4$ with mean $\Afit = 7.954 \approx 8 = 2d$
and standard deviation $3.7 \times 10^{-3}$ -- lies strictly
outside the formal baseline of \Cref{thm:baseline} at $d \ge 3$.
Two non-mutually-exclusive structural mechanisms are consistent
with this gap:
\begin{enumerate}[label=\textup{(\roman*$'$)},leftmargin=2.4em]
\item \emph{Borderline-locus mechanism.} The SIARC stratum is
positioned at the \emph{borderline} locus $\dega = 2\degb$ of
\cite[\S5]{wimp1984computation} (the Wimp--Zeilberger normal
form is replaced by an anormal-case ansatz of fractional rank
$q = 2$).
\item \emph{Definitional mechanism.} The empirical $\Afit$ of
\cite{papanokechi_pcf2_v13} measures a ratio different from
$\mu_{\mathrm{dom}} - \mu_{\mathrm{sub}}$ (for example, a
Stokes-constant exponent on $\log Q_n$ rather than on
$n \log n$).
\end{enumerate}
Distinguishing \textup{(i$'$)} from \textup{(ii$'$)} is
the open content of any follow-on phase
(\Cref{sec:open}, Q\ref{q:mechanism}).
\end{proposition}

\Cref{thm:baseline} is proven in \Cref{sec:proof} by explicit
case analysis on the three conventions across $d \in [2, 8]$;
the computation reproduces the symbolic derivation of
\cite[Phase~A]{siarc_t1_phase2} in self-contained form.
\Cref{prop:gap-framing} is supported in \Cref{sec:gap} by
direct comparison of the formal baseline with the
\cite{papanokechi_pcf2_v13} empirical record and a re-statement
of the borderline-locus condition of
\cite[\S5]{wimp1984computation}.

The exponent $A = 2d$ at $d \ge 3$ -- equivalently, the
\emph{closure} of the gap of \Cref{prop:gap-framing} via
mechanism \textup{(i$'$)} or \textup{(ii$'$)} -- is recorded
in this note as Conjecture B4 of \cite{papanokechi_pcf2_v13},
restated in \Cref{conj:b4} below. We do not claim to prove or
partially prove Conjecture B4: the formal-level baseline
\Cref{thm:baseline} narrows the structural framing of the gap
to the two mechanisms above and is consistent with -- but does
not establish -- $A = 2d$.

\begin{conjecture}[B4 of \cite{papanokechi_pcf2_v13}]
\label{conj:b4}
For the SIARC PCF Wallis recurrence \eqref{eq:wallis} at
$d \ge 3$, the empirical exponent $\Afit$ extracted from a
high-precision tail-window fit of $\log|p_n / p_{n-1}|$ takes
the value $A = 2d$ exactly, in the limit $n \to \infty$.
\end{conjecture}

\section{The Wimp--Zeilberger normal-case ansatz and its
         balance analysis}
\label{sec:wz-balance}

This section reproduces, in self-contained form, the
symbolic balance analysis of \cite[Phase~A]{siarc_t1_phase2}
that yields the $\mu$-pair \eqref{eq:mu-pair} and so the
band \eqref{eq:Anaive-band}.

\subsection{Substitution and the leading-order edge equation}
\label{sec:wz-edge}
Substitute the ansatz \eqref{eq:wz-ansatz} into the Wallis
recurrence \eqref{eq:wallis} and divide by $y(n) \ne 0$; using
the Stirling-type identity
$\Gamma(n-k)/\Gamma(n) \sim n^{-k}\,(1 + O(1/n))$ for fixed
$k \ge 1$ and the formal expansion
$(1 - 1/n)^{\rho} = 1 + O(1/n)$, the leading order in $n$ is
\begin{equation}
\label{eq:edge}
  1 \;-\; \frac{c_b\,n^{\degb - \mu}}{\gamma}
  \;-\; \frac{c_a\,n^{\dega - 2\mu}}{\gamma^{2}} \;=\; 0.
\end{equation}
Equation \eqref{eq:edge} is the leading-order edge equation
of the Newton-polygon balance for \eqref{eq:wallis} on the
ansatz \eqref{eq:wz-ansatz}; the three monomials $1$,
$c_b\,n^{\degb - \mu}/\gamma$, $c_a\,n^{\dega - 2\mu}/\gamma^{2}$
correspond to the three terms $p_n$, $b_n\,p_{n-1}$,
$a_n\,p_{n-2}$ of \eqref{eq:wallis}.

\subsection{Three balance possibilities}
\label{sec:wz-three-balances}
The leading-order edge equation \eqref{eq:edge} admits three
possible \emph{Newton-polygon balances} -- pairings of two of
the three monomials at the same order in $n$, with the third
monomial subleading:

\begin{itemize}[leftmargin=2.2em]
\item[\textbf{(I)}] \emph{$1 \leftrightarrow c_b/\gamma$ balance.}
   $\degb - \mu = 0$, equivalently $\mu = \degb = d$ and
   $\gamma = c_b$. The third term
   $c_a\,n^{\dega - 2\mu}/\gamma^{2}$ is subleading whenever
   $\dega < 2\degb$.
\item[\textbf{(II)}] \emph{$1 \leftrightarrow c_a/\gamma^{2}$ balance.}
   $\dega - 2\mu = 0$, equivalently $\mu = \dega / 2$ and
   $\gamma^{2} = c_a$. The third term
   $c_b\,n^{\degb - \mu}/\gamma$ is subleading whenever
   $\degb < \dega / 2$, i.e., $\dega > 2\degb$.
\item[\textbf{(III)}] \emph{$c_b/\gamma \leftrightarrow c_a/\gamma^{2}$
   balance.} $\degb - \mu = \dega - 2\mu$, equivalently
   $\mu = \dega - \degb$ and $\gamma = -c_a / c_b$. The third
   term $1$ is subleading whenever $\degb - \mu > 0$, i.e.,
   $\mu < \degb$, i.e., $\dega < 2\degb$.
\end{itemize}

For the SIARC PCF stratum at $d \ge 2$ in any of the three
conventions ($\dega \in \{d-1, d, d+1\}$), the inequality
$\dega < 2\degb = 2d$ is satisfied (recalling $d \ge 2$, hence
$\dega \le d+1 \le 2d - 1 < 2d$). Therefore balance \textbf{(I)}
and balance \textbf{(III)} are simultaneously consistent, and
balance \textbf{(II)} is excluded. We obtain exactly two formal
solutions at the level of the leading-order edge equation
\eqref{eq:edge}:
\begin{equation}
\label{eq:dom-sub}
  \mu_{\mathrm{dom}} \;=\; d
  \quad (\text{from (I)}),
  \qquad
  \mu_{\mathrm{sub}} \;=\; \dega - d
  \quad (\text{from (III)}).
\end{equation}
The dominant solution $y_{\mathrm{dom}}$ corresponds to balance
\textbf{(I)} and carries the \emph{larger} $\mu$, hence the
\emph{larger} formal $\Gamma(n)^{\mu}$ growth rate; the
subdominant solution $y_{\mathrm{sub}}$ corresponds to balance
\textbf{(III)} and carries the \emph{smaller} $\mu$. The
inequality $\mu_{\mathrm{dom}} > \mu_{\mathrm{sub}}$ holds
strictly: $d - (\dega - d) = 2d - \dega \ge 1 > 0$ since
$\dega \le d + 1 \le 2d - 1$.

\subsection{The per-convention table}
\label{sec:wz-table}
Substituting the three SIARC conventions into \eqref{eq:dom-sub}
and \eqref{eq:Anaive-def}--\eqref{eq:Anaive-band} yields the
table:

\begin{center}
\begin{tabular}{c c c c c c}
\toprule
Convention & $\dega$ & $\mu_{\mathrm{dom}}$ & $\mu_{\mathrm{sub}}$
  & $\Anaive = 2d - \dega$ & $\Anaive$ as $f(d)$\\
\midrule
$\alpha$-direction & $d-1$ & $d$ & $-1$ & $d + 1$ & $d+1$\\
symmetric          & $d$   & $d$ & $0$  & $d$     & $d$\\
$\delta$-direction & $d+1$ & $d$ & $1$  & $d - 1$ & $d-1$\\
\bottomrule
\end{tabular}
\end{center}

The three columns labelled $\Anaive$ exhibit the band
$\Anaive \in \{d-1,\,d,\,d+1\}$ uniformly in $d$; the
band's width is $2$ (from $d-1$ to $d+1$), independent of $d$.
The largest entry in the band, $\Anaive = d+1$, is the
\emph{$\alpha$-direction} value, and is the only entry that
recovers an integer in the range of the empirical $\Afit$
for $d = 2$ (see \Cref{sec:proof-d2}); for $d \ge 3$, all
three entries lie strictly below the empirical $\Afit \approx 2d$
of \cite{papanokechi_pcf2_v13}.

\subsection{Phase~A AEAL anchor}
\label{sec:wz-aeal}
The symbolic derivation above is implemented in the Phase~A
script \texttt{phase\_a\_symbolic\_derivation.py} of the
T1-BIRKHOFF-PHASE2-LIFT-LOWER bridge session
\cite{siarc_t1_phase2}, locked at SHA-$256$
\texttt{2fc6e392\dots f16248} (full hash:
\texttt{2fc6e39267768791912cc53aa59bc231c40483a1a93e92104b13f07809f16248}).
The script's run-output table at $d \in \{2, 3, 4\}$ for the three
conventions reproduces the table of \Cref{sec:wz-table} verbatim;
the AEAL claim entries \texttt{PA\_DEG2\_ALPHA\_BASELINE},
\texttt{PA\_DEG3\_ALPHA\_BASELINE}, \texttt{PA\_DEG4\_ALPHA\_BASELINE}
of the bridge \texttt{claims.jsonl} carry the per-degree
output hashes.

\section{Proof of the formal-level baseline}
\label{sec:proof}

\Cref{thm:baseline} is proven by combining the symbolic balance
analysis of \Cref{sec:wz-balance} (which establishes the
$\mu$-pair \eqref{eq:dom-sub} and the band
\eqref{eq:Anaive-band} for the leading-order edge equation
\eqref{eq:edge}) with the per-degree case analysis below.

\subsection{Per-degree confirmation: $d \in [2, 8]$}
\label{sec:proof-d28}
The Phase~A symbolic derivation
\cite[Phase~A.1--A.4]{siarc_t1_phase2} executes the substitution
of \Cref{sec:wz-edge} for $d \in \{2, 3, 4\}$ in the three
SIARC conventions. The Phase~B extended sweep
\texttt{phase\_b\_extended\_sweep.py} of
\cite[Phase~B]{siarc_t1_phase2} (SHA-$256$
\texttt{39e98db607a5331686cf4ba4890bc50abc1232e8b533b80eeb2d4636c8e10175})
extends the same substitution to $d \in [3, 8]$; in every
$(d, \text{convention})$ stratum sampled, the dominant balance
is \textbf{(I)} and the subdominant balance is \textbf{(III)}
of \Cref{sec:wz-three-balances}, and the band
$\Anaive \in \{d-1, d, d+1\}$ is realised exactly. The Phase~B
AEAL claim \texttt{PB\_SWEEP\_D5\_TO\_D8} of the bridge
\texttt{claims.jsonl} carries the corresponding output hash.

The Phase~B sweep also \emph{records} that the borderline
condition $\dega = 2\degb$ of
\cite[\S5]{wimp1984computation} is \emph{never} met by any of
the three conventions for $d \in [3, 8]$ (since $\dega \le d+1 <
2d$ strictly), so that the SIARC stratum sits in the
\emph{normal} case of the Wimp--Zeilberger formalism throughout
this range.

This per-degree case analysis closes the proof of
\Cref{thm:baseline}.

\subsection{Verification at $d = 2$, $\alpha$-direction}
\label{sec:proof-d2}
Specialising \Cref{thm:baseline} to $d = 2$ in the
$\alpha$-direction convention ($\dega = d - 1 = 1$) yields
$\Anaive = d + 1 = 3$. This formal-level prediction is the
exponent recovered, at $80$ algebraic digits, by the
PCF-1 v1.3 \S6 Theorem 5 \emph{lower branch}:

\begin{quote}
\itshape
At $d = 2$, the $\alpha$-direction Wallis stratum admits two
exponent strata, $A \in \{3, 4\}$. The lower stratum $A = 3$ is
realised by the QL01--QL26 family of $26$ representatives, with
each representative's $A$-value verified at $\dps = 80$ across
the high-precision tail-window fit (\cite[\S6, Theorem~5,
lower branch]{papanokechi_pcf1_v13}).
\end{quote}

The Phase~A claim \texttt{PA\_DEG2\_ALPHA\_BASELINE} of
\cite{siarc_t1_phase2} records that the symbolic derivation
of \Cref{sec:wz-balance} at $d = 2$, $\alpha$-direction
recovers $\Anaive = 3$ verbatim, matching the QL01--QL26 lower
branch as verified at $\dps = 80$ in \cite{papanokechi_pcf1_v13}.

The $d = 2$ \emph{upper} branch $A = 4 = 2d$, realised by the
V$_{\mathrm{quad}}$ representative of
\cite[\S6, Theorem~5, upper branch]{papanokechi_pcf1_v13},
is \emph{not} recovered by the formal-level baseline of
\Cref{thm:baseline}: the band $\Anaive \in \{1, 2, 3\}$ at
$d = 2$ does not contain $4$. The V$_{\mathrm{quad}}$ upper
branch is consistent with the borderline-locus mechanism
\textup{(i$'$)} of \Cref{prop:gap-framing} -- the closure of
which is the open content of \Cref{sec:open} Q\ref{q:mechanism}.

\section{Structural framing of the gap to $A = 2d$ at $d \ge 3$}
\label{sec:gap}

\Cref{thm:baseline} establishes that the formal-level
baseline $\Anaive$ obtained from the Wimp--Zeilberger normal-case
ansatz lies in the band $\{d-1, d, d+1\}$ across the three SIARC
conventions, uniformly for $d \in [2, 8]$. The empirical
record of PCF-2 v1.3 R1.1 + R1.3 + Q1 \cite{papanokechi_pcf2_v13}
exhibits an $\Afit \approx 2d$ at $d = 3$ and $d = 4$ -- a value
that lies strictly outside this band for $d \ge 3$. This
section frames the gap structurally; it does \emph{not} close
the gap.

\subsection{The PCF-2 v1.3 empirical record}
\label{sec:gap-pcf2}
The empirical record carried by \cite{papanokechi_pcf2_v13}
R1.1 + R1.3 + Q1 is, verbatim:

\begin{itemize}[leftmargin=2.2em]
\item $50/50$ cubic representatives at $d=3$ with
  $\Afit \approx 6 = 2d$ in the high-precision tail-window
  fit (R1.1 of \cite{papanokechi_pcf2_v13}).
\item $60/60$ quartic representatives at $d=4$ with mean
  $\Afit = 7.954 \approx 8 = 2d$ and standard deviation
  $\sigma = 3.7 \times 10^{-3}$ (R1.3 + Q1 of
  \cite{papanokechi_pcf2_v13}).
\end{itemize}

The fit is performed at $\dps = 80$ in mpmath, on a tail
window of length $O(10^3)$ to $O(10^4)$, with the $A$-value
extracted by linear regression on
$\log|p_n / p_{n-1}|$ against $n \log n$ (cf.\
\cite[\S\S R1, Q1]{papanokechi_pcf2_v13} for the exact
window-length and fit-protocol). The Phase~A baseline of
\Cref{thm:baseline} predicts $\Anaive \in \{d-1, d, d+1\}$
in the same conventions, hence $\Anaive < 2d$ for $d \ge 3$;
the record is consistent with the empirical lift but is not
explained at the formal level.

\subsection{The Wimp--Zeilberger borderline-case condition}
\label{sec:gap-wz-border}
\cite[\S5]{wimp1984computation} distinguishes a \emph{normal}
case ($\dega < 2\degb$) from a \emph{borderline} case
($\dega = 2\degb$) of the second-order linear difference
recurrence. In the normal case, the Wimp--Zeilberger ansatz
\eqref{eq:wz-ansatz} of integer formal rank $p = 1$ produces
exactly the two formal solutions $y_{\mathrm{dom}}$ and
$y_{\mathrm{sub}}$ of \Cref{sec:wz-three-balances}. In the
borderline case, the two formal solutions \emph{coalesce} at
the leading symbol level (both have $\mu = \degb$), and the
asymptotic split is governed instead by sub-leading
sub-exponential corrections of the form $\exp(\pm B\sqrt{n})$;
the appropriate ansatz is anormal of fractional rank $q = 2$
(cf.\ \cite[\S\S5--6]{wimp1984computation}).

The Phase~B sweep \cite[Phase~B]{siarc_t1_phase2} records that
the SIARC stratum at $d \ge 2$ in any of the three SIARC
conventions has $\dega \in \{d-1, d, d+1\}$, hence
$\dega < 2\degb = 2d$ \emph{strictly} for $d \ge 2$. The SIARC
stratum is therefore \emph{in} the normal case of
\cite[\S5]{wimp1984computation}, and the formal-level
baseline of \Cref{thm:baseline} is the formal exponent
predicted by the normal-case ansatz.

\subsection{Two structural mechanisms consistent with the gap}
\label{sec:gap-mechanisms}

The gap between $\Anaive \in \{d-1, d, d+1\}$ (formal baseline)
and $\Afit \approx 2d$ (empirical record at $d = 3, 4$) at
$d \ge 3$ is consistent with two non-mutually-exclusive
structural mechanisms:

\begin{itemize}[leftmargin=2.2em]
\item[\textbf{(i$'$)}] \emph{Borderline-locus mechanism.}
   The SIARC stratum may exhibit, at $d \ge 3$, a degeneracy
   of the leading polynomial coefficient $c_b$ on a particular
   sub-locus of the parameter space, in such a way that the
   effective $\dega$ at that sub-locus equals $2\degb$. The
   Wimp--Zeilberger normal-case ansatz then fails on the
   sub-locus and is replaced by an anormal-case ansatz of
   fractional rank $q = 2$ (carrying $\exp(\pm B\sqrt{n})$
   sub-leading factors). The SIARC stratum's empirical
   $\Afit \approx 2d$ would, on this mechanism, reflect the
   $\mu$-pair of the anormal-case ansatz and not the
   $\mu$-pair \eqref{eq:dom-sub} of the normal case.
\item[\textbf{(ii$'$)}] \emph{Definitional mechanism.}
   The empirical $\Afit$ extracted by the PCF-2 v1.3
   high-precision tail-window fit may measure a ratio
   different from $\mu_{\mathrm{dom}} - \mu_{\mathrm{sub}}$.
   For example, $\Afit$ may be the Stokes-constant exponent
   on $\log|p_n|$ -- a slightly different combination of the
   formal data of \eqref{eq:wz-ansatz} -- rather than the
   exponent on $n \log n$. On this mechanism, the formal
   baseline of \Cref{thm:baseline} and the empirical
   $\Afit \approx 2d$ would not be in direct tension; the
   gap would instead reflect a translation between two
   distinct asymptotic invariants.
\end{itemize}

These two mechanisms are not mutually exclusive: a sub-locus
identification of the form (i$'$) may coexist with a
definitional translation of the form (ii$'$). The argument
supporting \Cref{prop:gap-framing} consists of the two
observations above; the closure of the gap (i.e., the
identification of which mechanism is realised) is the open
content of \Cref{sec:open} Q\ref{q:mechanism}.

\subsection{Forbidden-verb hygiene for this section}
\label{sec:gap-hygiene}
The structural framing of \Cref{sec:gap-mechanisms} does
not assert that either mechanism (i$'$) or (ii$'$) is
realised. The wording is \emph{consistent with}, \emph{frames
the gap}, \emph{narrows the gap}, \emph{is consistent with};
the wording \emph{shows / confirms / proves / demonstrates /
establishes / verifies} is reserved for the formal-level
\Cref{thm:baseline} and the literature-citation
$d=2$ verification of \Cref{sec:proof-d2}, and is not
extended to $A = 2d$ or to Conjecture B4
(\Cref{conj:b4}). This hygiene is checked, sentence by
sentence, in the \texttt{rubber\_duck\_critique.md} of
this session.

\section{Connection to Birkhoff--Trjitzinsky 1933}
\label{sec:bt}

The Birkhoff--Trjitzinsky 1933 paper
\cite{birkhoff_trjitzinsky_1933} provides the
\emph{existence and factorization} machinery for the formal
solutions of a difference equation $L_n(y) = 0$ with
coefficients of the kind specified in \cite[\S1]{birkhoff_trjitzinsky_1933}.
This section records that the machinery applies to the SIARC
stratum uniformly in $d$ at the formal level and the sectorial
realization level. It does \emph{not} record that the machinery
identifies $A = 2d$ for the SIARC stratum specifically; that
identification is the open content of any follow-on phase
(\Cref{sec:open}).

\subsection{The B--T 1933 \S\S7--9 machinery}
The Phase~C verification \cite[Phase~C]{siarc_t1_phase2}
extracts verbatim from the SHA-verified OCR text dump of
\cite{birkhoff_trjitzinsky_1933} (PDF SHA-$256$
\texttt{dcd7e3c6\dots d68fe6}) the following three results:

\begin{itemize}[leftmargin=2.2em]
\item \textbf{\S7 Theorem II} (factorization preservation):
  if $L_n = L_{n-r}\,L_r$ with $L_{n-r}, L_r$ completely proper
  in a region, the product $L_n$ is completely proper in the
  same region.
\item \textbf{\S9 Fundamental Theorem} (existence):
  every $L_n(y) = 0$ with coefficients of the kind specified
  in \cite[\S1]{birkhoff_trjitzinsky_1933} and known in the
  complete neighbourhood of infinity is completely proper in
  each of the several quadrants associated with the equation.
\item \textbf{\S11 Theorem III} (converse / sectorial
  realization): given a linearly independent set of formal
  series $e^{Q_j(x)}\,s_j(x)$ of the same general description
  as might occur in connection with a difference equation of
  order $n$, there exists a difference equation $L_n(y) = 0$
  of \cite[\S1]{birkhoff_trjitzinsky_1933}-kind possessing the
  series as fundamental solutions.
\end{itemize}

\subsection{Application to the SIARC stratum}
The SIARC PCF Wallis recurrence \eqref{eq:wallis} satisfies
the \cite[\S1]{birkhoff_trjitzinsky_1933}-kind hypothesis: its
coefficients $a_n, b_n$ are polynomial in $n$ (and in
particular asymptotically rational in $n$ at $n = \infty$),
hence of the appropriate Gevrey-$1$ class. The order is
$n = 2$ in the language of \cite{birkhoff_trjitzinsky_1933}
(an order-$2$ linear difference operator); the polynomial
degree $d$ of the coefficient $b_n$ enters only at the level
of the leading symbol of the operator, not the
\cite[\S1]{birkhoff_trjitzinsky_1933}-kind hypothesis itself.

The Phase~C four sub-gates of \cite[Phase~C.2]{siarc_t1_phase2}
(non-resonance match, non-degeneracy match, $d$-range
applicability, $\mu$-units consistency per the upstream A-01
verdict) all pass for the SIARC stratum at the
\emph{normal-case} locus $\dega < 2\degb$ -- with a
\emph{caveat} that the \S7 Theorem~II ``point of division''
condition $\Re Q_{\mathrm{dom}}(x) > \Re Q_{\mathrm{sub}}(x)$
\emph{degenerates} on the borderline locus $\dega = 2\degb$
(both formal exponents $\mu_{\mathrm{dom}}, \mu_{\mathrm{sub}}$
coalesce at $\degb$, and the strict inequality on the real
parts of the exponential dominants is replaced by an equality
at leading order). The borderline-locus degeneracy is the
literature-side counterpart of the borderline-locus mechanism
\textup{(i$'$)} of \Cref{prop:gap-framing}.

\subsection{Scope of the application}
The application records that the formal-level analysis of
\Cref{thm:baseline} carries through the
\cite{birkhoff_trjitzinsky_1933} \S\S7--9 existence and
factorization machinery uniformly in $d \in [2, 8]$. It
\emph{does not} record that \cite{birkhoff_trjitzinsky_1933}
identifies $A = 2d$ for the SIARC stratum specifically; the
\S9 Fundamental Theorem is universal in $d$ but pin-points
the formal-level $\mu$-pair only via the
\cite[\S1]{birkhoff_trjitzinsky_1933}-kind hypothesis,
which gives the band $\Anaive \in \{d-1, d, d+1\}$ already
recorded in \Cref{thm:baseline}, not the lift to
$A = 2d$. The lift is open in this note and is the open
content of \Cref{sec:open} Q\ref{q:mechanism},
Q\ref{q:borderline-Q}, and Q\ref{q:sectorial}.

\section{Open questions and follow-on phases}
\label{sec:open}

\begin{problem}[Mechanism identification]
\label{q:mechanism}
Which of the two mechanisms (i$'$) borderline-locus and
(ii$'$) definitional of \Cref{prop:gap-framing} is realised
on the SIARC PCF Wallis stratum at $d \ge 3$? The
\cite{papanokechi_pcf2_v13} R1.1 + R1.3 + Q1 record is
consistent with either; distinguishing them requires
either (a) symbolic derivation of the $\sqrt{n}$ sub-leading
correction predicted by mechanism (i$'$) and an empirical
test on the high-precision tail of $\log|p_n / p_{n-1}|$, or
(b) a careful re-statement of the \cite{papanokechi_pcf2_v13}
fit-protocol to support or rule out the
$\Afit = \mu_{\mathrm{dom}} - \mu_{\mathrm{sub}}$ identification
under mechanism (ii$'$).
\end{problem}

\begin{problem}[Borderline-case algebraic ansatz]
\label{q:borderline-Q}
Under mechanism (i$'$), can the borderline-case algebraic
ansatz
$Q_j(x) = \pm B\,x^{1/2} + (\text{lower order})$
of \cite[\S1]{birkhoff_trjitzinsky_1933} be made explicit at
the SIARC stratum, with $B$ identified in closed form in terms
of $c_a$ and $c_b$? The expected closed form is
$B = \sqrt{c_a}$ at the formal level, by analogy with the
$d = 2$ V$_{\mathrm{quad}}$ upper-branch transition of
\cite[\S6]{papanokechi_pcf1_v13}; an explicit derivation is
open in this note.
\end{problem}

\begin{problem}[Definitional mechanism]
\label{q:Afit-def}
Under mechanism (ii$'$), what is the precise mathematical
definition of $\Afit$ extracted by the
\cite{papanokechi_pcf2_v13} high-precision tail-window fit?
The fit is performed on $\log|p_n / p_{n-1}|$ at $\dps = 80$;
the linear-regression target is $n \log n$. Whether the
extracted slope is $\mu_{\mathrm{dom}} - \mu_{\mathrm{sub}}$
or instead a Stokes-constant exponent (or a combination of
the two) requires a fit-protocol audit.
\end{problem}

\begin{problem}[Formal-to-analytic sectorial upgrade]
\label{q:sectorial}
The formal-level baseline of \Cref{thm:baseline} is purely
formal: it produces the two formal solutions
$y_{\mathrm{dom}}, y_{\mathrm{sub}}$ as elements of
$\Gamma(n)^{\mu}\,\gamma^{n}\,n^{\rho}\,\mathbb{C}[[1/n]]$,
without an analytic-side guarantee that the corresponding
formal Borel transforms converge in a sector of opening
$> \pi/(2d)$ (the canonical sectorial-asymptotic existence
condition for irregular singularities of formal rank $p = 1$).
The formal-to-analytic sectorial upgrade is the content of
the Wasow~\S X.3 (Theorem~11.1) \cite{wasow1965asymptotic} /
Adams~1928 \cite{adams1928analytic} /
Turrittin~1955 \cite{turrittin1955convergent} /
Immink~1984 \cite{immink1984asymptotics} framework, with the
modern Borel-Laplace radius identification stated in
\cite[ch.~5]{costin2008asymptotics} and the formal-side
existence anchored in \cite[\S2]{birkhoff1930formal}
(see \cite[Phase~D]{siarc_t1_phase2} for the literature
substrate).
At the time of preparation of this note, the
Wasow~\S X.3 OCR-readable PDF is unavailable on disk
(image-only) and \cite{adams1928analytic} has not been
acquired (NIA per the upstream A-01 verdict of
\cite{siarc_t1_phase2}).
The sectorial upgrade is open and is the natural T1 Phase~3
follow-on.
\end{problem}

\section*{AI Disclosure}

During the preparation of this work the author used GitHub
Copilot (Microsoft) and Anthropic Claude to assist with code
generation, numerical computations, and manuscript drafting.
After using these tools, the author reviewed and edited the
content as needed and takes full responsibility for the
content of this article. AI assistants are not authors. All
theorem statements (\Cref{thm:baseline}), proposition
statements (\Cref{prop:gap-framing}), and conjecture
statements (\Cref{conj:b4}) remain the human author's
responsibility, and every numerical or symbolic-derivation
fact in this note traces to an AEAL entry in
\texttt{claims.jsonl} with a SHA-$256$ output hash on the
SIARC bridge under \texttt{sessions/2026-05-03/T1-BIRKHOFF-PHASE2-LIFT-LOWER/}
(the Phase~2 substrate, commit \texttt{37c939f}) and
\texttt{sessions/2026-05-06/T1-PHASE2-BASELINE-NOTE/} (this
version). The Birkhoff--Trjitzinsky~1933~\S\S7--9 verbatim
extracts of \Cref{sec:bt} are anchored by the SHA-verified PDF
\texttt{dcd7e3c6\dots d68fe6} of
\cite{birkhoff_trjitzinsky_1933} at the runbook-canonical
literature directory
(\texttt{tex/submitted/control center/literature/g3b\_2026-05-03/}).

\medskip
\noindent\emph{Related cross-cites for the Zenodo
related-identifiers block.}
The present note's Zenodo deposit lists the following six
SIARC Zenodo records as related identifiers (cite-block):
PCF-1 v1.3 \cite{papanokechi_pcf1_v13},
PCF-2 v1.3 \cite{papanokechi_pcf2_v13},
Channel Theory v1.3 \cite{siarc_channel_theory_v13},
T2B v3 \cite{siarc_t2b_v3},
SIARC umbrella v2.0 \cite{siarc_umbrella_v2}, and
D2-NOTE v2.1 \cite{siarc_d2_note_v21}.

\appendix

\section{Bridge session provenance}
\label{app:provenance}

The following SIARC bridge sessions are the AEAL anchors for
the per-theorem claims in this note. Each is committed to the
public \texttt{siarc-relay-bridge} GitHub repository at
\href{https://github.com/papanokechi/siarc-relay-bridge}{\nolinkurl{github.com/papanokechi/siarc-relay-bridge}}.

\begin{itemize}[leftmargin=2.2em]
\item \texttt{sessions/2026-05-03/T1-BIRKHOFF-PHASE2-LIFT-LOWER/}
  (commit \texttt{37c939f}, 2026-05-04). Phase~2 of T1; verdict
  \texttt{UPGRADE\_PARTIAL\_FORMAL\_LEVEL}. The substrate of this
  note. Phase~A symbolic derivation
  \texttt{phase\_a\_symbolic\_derivation.py} SHA-$256$
  \texttt{2fc6e392\dots f16248}; Phase~B extended sweep
  \texttt{phase\_b\_extended\_sweep.py} SHA-$256$
  \texttt{39e98db6\dots 8c10175}; Phase~C literature
  verification \texttt{phase\_c\_literature\_verification.md}
  SHA-$256$ \texttt{c6cc2b1a\dots 808ae93}; Phase~D verdict
  \texttt{phase\_d\_verdict.md} SHA-$256$
  \texttt{7145a00d\dots 2c65c5}; Phase~2 \texttt{claims.jsonl}
  carries 16 AEAL entries.
\item \texttt{sessions/2026-05-03/T1-A01-NORMALIZATION-RESOLUTION/}.
  Upstream A-01 normalization-resolution verdict, used in
  \Cref{sec:bt} to confirm $\mu$-units consistency across
  Wasow / Birkhoff / Birkhoff--Trjitzinsky / Adams normalisation
  conventions; provides the SHA-verified
  \cite{birkhoff_trjitzinsky_1933} OCR text dump
  \texttt{\_bt1933.txt} that supplies the verbatim extracts
  of the Phase~C verification.
\item \texttt{sessions/2026-05-06/T1-PHASE2-BASELINE-NOTE/} (this
  session). Drafting and Zenodo packaging of the present note.
  Contains \texttt{bt\_baseline\_note.tex},
  \texttt{bt\_baseline\_note.pdf},
  \texttt{annotated\_bibliography.bib},
  \texttt{zenodo\_description\_bt\_baseline\_note.txt},
  \texttt{zenodo\_upload\_bt\_baseline\_note\_runbook.md},
  \texttt{claims.jsonl} ($\ge 10$ AEAL entries),
  \texttt{rubber\_duck\_critique.md}, \texttt{verdict.md},
  \texttt{halt\_log.json}, \texttt{discrepancy\_log.json},
  \texttt{unexpected\_finds.json}, and \texttt{handoff.md}.
\end{itemize}

\bibliographystyle{plain}
\bibliography{annotated_bibliography}

\end{document}
